团队在比两个语言模型。A 的困惑度 12.4，B 是 11.8。结论是 B 更好——这没错，但多数人说不清好在哪：这 0.6 意味着什么，值多少工程投入？

信息论的回答很具体。困惑度是交叉熵的指数，而交叉熵可以拆成两段：一段是数据本身的熵，谁也压不掉；另一段是模型偏离真实分布的代价，单位是比特。所以两个模型的差距不是一个抽象分数，它是\emph{每个词平均多花的比特}——同时也是压缩这份语料时多占的存储。

换个小一点的例子把这笔账算完。一枚正面概率 \(0.9\) 的硬币，每次抛掷的熵是 \(h_2(0.9)\approx0.469\) 比特。若编码者误以为它公平，每次固定花 \(1\) 比特，多付的 \(0.531\) 比特恰好就是这两个分布之间的 KL 散度。模型错在哪里、错了多少，被换算成了一个可以计价的数。

这不是巧合，而是这套语言的全部用途：\textbf{把``不确定''和``不知道''变成能加减、能设界、能比较的量。}一旦不确定性有了单位，原本只能定性争论的问题就有了确定答案——这份数据最多压到多小、这条信道最多传多快、这组特征里究竟有没有关于标签的信息。

\begin{center}
\begin{tikzpicture}[x=1cm,y=1cm,font=\small,>=stealth]
  \draw[black!45,fill=blue!9] (0,0.6) rectangle (3.5,2.0);
  \node[align=center] at (1.75,1.3) {不确定性\\写成比特};
  \draw[black!45,fill=green!8] (4.8,2.55) rectangle (9.6,3.6);
  \node[align=center] at (7.2,3.07) {压缩：平均码长 \(\ge H\)};
  \draw[black!45,fill=green!8] (4.8,0.78) rectangle (9.6,1.83);
  \node[align=center] at (7.2,1.3) {通信：可靠速率 \(\le C\)};
  \draw[black!45,fill=green!8] (4.8,-0.99) rectangle (9.6,0.06);
  \node[align=center] at (7.2,-0.47) {学习：错误率 \(\ge\) Fano 下界};
  \draw[->,black!55] (3.55,1.45) -- (4.75,3.05);
  \draw[->,black!55] (3.55,1.30) -- (4.75,1.30);
  \draw[->,black!55] (3.55,1.15) -- (4.75,-0.45);
  \draw[black!30,dashed] (9.6,1.3) -- (10.35,1.3);
  \node[align=left,anchor=west,font=\footnotesize,black!70] at (10.45,1.3)
    {三条边界共用一个量纲，\\于是``还能不能再好一点''\\成了可回答的问题};
\end{tikzpicture}

\emph{压缩、通信、学习看着是三件事，给出的却是同一种断言：某个量存在硬下界，越过它的努力必然徒劳。知道边界在哪，比多试几次省事得多。}
\end{center}

\subsection{五条贯穿始终的判断}

\textbf{熵是``平均要问多少个是非题''，不是抽象的混乱程度。}这个操作性解释让比特成为真正的单位，也立刻把熵与编码接上：最优码长就是自信息 \(-\log_2p\)，而 Kraft 不等式说一组码长可行当且仅当 \(\sum2^{-\ell_i}\le1\)。分配码长与分配概率，是同一件事的两种写法。

\textbf{熵只依赖概率，不依赖取值。}把数值放大一千倍，方差涨一百万倍，熵纹丝不动。这是长处也是短处：它度量不确定性，不度量大小。连带的后果是微分熵不能单独使用——换个单位它就变，所以有意义的只有差值，互信息与 KL 在坐标变换下不变，微分熵不变不了。

\textbf{KL 是``用错模型的代价''，不是分布之间的距离。}上面那 \(0.531\) 比特就是它。这个读法立刻解释了不对称性，也解释了两个方向为何导致不同的失败：一个逼着模型覆盖全部支撑，症状是生成模糊；一个允许模型缩到单峰，症状是只会产出同一种东西。看到这两类症状，先回头看目标函数用的是哪个方向。

\textbf{加工不会增加信息。}数据处理不等式说 \(X\to Y\to Z\) 时 \(I(X;Z)\le I(X;Y)\)：任何确定性变换都造不出关于标签的新信息，能改善的只是可提取性。反过来更有用——如果处理之后互信息\emph{变大}了，那不是发现，是 Markov 条件被破坏，通常意味着信息泄漏。

\textbf{边界是硬的，而反面结论同样值钱。}容量之下可以做到任意可靠，容量之上无论怎样改进编码都压不下错误率；Fano 不等式说特征没覆盖住的那部分标签不确定性，会转化成错误率的下界。这条判断区分了``模型还不够好''与``信息本来就不够''——两种情况需要的下一步完全不同，而实践中它们经常被混为一谈。

\subsection{六章怎么安排}

\hyperref[ch:054]{``信息与熵：怎样把``不确定''变成可计算的量''}是地基。其中微分熵那一篇容易被跳过，但它讲的正是连续型变量为什么不能直接比较``信息量''，这个坑很多人踩过。

\hyperref[ch:055]{``联合信息：一个变量能告诉我们另一个变量多少''}与\hyperref[ch:056]{``分布之间的差异：从错误编码到散度''}使用频率最高。前者给出比相关系数更可靠的依赖度量，后者解释了交叉熵损失、变分推断、生成模型目标函数的共同来源。做机器学习的读者若只读两章，就读这两章。

\hyperref[ch:057]{``无损信源编码：熵为什么是压缩极限''}与\hyperref[ch:058]{``有噪信道编码：噪声中为什么仍能可靠通信''}是两条经典定理线。工程细节可以略读，两条定理本身值得记住——尤其后者那个反直觉结论：信道会出错，却仍能做到任意可靠，让出的只是速率，而不是无限增加的冗余。

\hyperref[ch:059]{``信息与学习：哪些数据值得看，哪些表示值得保留''}把前五章接到机器学习上。这一章里有定论的部分（交叉熵与 KL 的等价、最小描述长度、Fano 界）和仍有争议的部分（信息瓶颈在确定性网络中的解释）已经分开写，读的时候留意这个区分——把有争议的当定论转述，是这个领域最常见的失真方式。

读这一部分时值得反复问一句：手上这个量的单位是什么？答案若是比特，它就能与另一个比特量相加、相减、设界；答案若是``某种分数''，那多半还没真正进入信息论的语言。
